\section{Further certified families}
\label{app:research-extensions}
\label{app:projective-cubes}
\label{app:research-provenance-corrections}
Beyond the three case studies of Section~\ref{sec:mathematics}, further witnesses from the research collection were extended to parameter families and formalized. Table~\ref{tab:certified-research-extensions} summarizes them, and the paragraphs that follow give each result's exact scope, its certificate, and where its witness came from. The prime-blowup and rainbow-Schur families are stated and proved in Appendix~\ref{app:structural-case-studies}. We developed these extensions with the help of AI tools. Separate adversarial model reviews checked their formal statements and their correspondence to the sources. They are mathematical results, not measurements of model discovery. None of the files uses \texttt{sorry} or \texttt{native\_decide}. \path{mathematics/index.json}, \path{mathematics/extensions/index.json}, and \path{mathematics/supporting/index.json} link the statements, sources, hashes, and build records, and \path{mathematics/discoveries/} records the source versions and discovery provenance. Mathematical priority is left unassigned throughout.

\begin{table}[ht]
\centering\small
\setlength{\tabcolsep}{4pt}
\begin{tabular}{@{}p{.16\linewidth}p{.28\linewidth}p{.31\linewidth}p{.16\linewidth}@{}}
\toprule
Case & Refuted statement & What the selected Lean proof establishes & Toolchain \\
\midrule
Projective cubes & Conjecture~4 of \citet{meng2026cubefree} & Cube-free sets above the conjectured density for $d=5,7,8$ in infinitely many cyclic groups; the exact density boundary of the layered construction & core Lean 4.30.0 and 4.34.0; boundary results Lean 4.30.0 with Mathlib \\
Supertoken graphs & $\dim F_2(K_n)=n-1$ \citep{baskoro2025supertoken} & Resolving sets of size at most $n-2$ for every $n\ge8$ & core Lean 4.30.0 and 4.34.0 \\
Binomial polynomials & second assertion of \citet{sun2025polynomials} & No coefficientwise $q$-log-convexity for every integer $m\ge4$ & core Lean 4.30.0 and 4.34.0 \\
Fake-Gaussian positivity & Conjecture~4 of \citet{gatzweilerkrattenthaler2025positivity} & An exact integer polynomial with a negative coefficient & Lean 4.30.0 and 4.34.0 \\
$q$-rious unimodality & unimodality assertion of \citet{warnaarzudilin2025qrious} & A non-unimodal coefficient sequence; the Landau hypothesis certified & Lean 4.30.0 and 4.34.0 \\
Permutation-grid rooks & necessity direction of Conjecture~3.10 of \citet{lewiswon2026permutationgrids} & A placement count congruent to 3 modulo 4 & Lean 4.30.0 and 4.34.0 \\
Toda numerators & Conjecture~7.3 of \citet{labelle2026toda} & A numerator that is not unimodal under any common-exponent regrading & Lean 4.30.0 and 4.34.0 \\
\bottomrule
\end{tabular}
\caption{Further certified families and refutations. Each row's scope, certificate file, provenance, and the human judgments on its dossier are given in the text; the prime-blowup and rainbow-Schur certificates are in Appendix~\ref{app:structural-case-studies}.}
\label{tab:certified-research-extensions}
\end{table}

\paragraph{Projective cubes.}
In the cyclic group $\mathbb Z_N$, the projective cube of a size-$d$ multiset $S=\{a_1,\ldots,a_d\}$ is $\Sigma^*S=\{\sum_{i\in I}a_i:\emptyset\ne I\subseteq[d]\}$, the set of all sums over nonempty index sets. Repeated and zero generators are allowed, and the source prints $\subset$ in its definition but its three-generator example includes $x+y+z$, so we include the full index set. A set is cube-free if it contains no such cube. Conjecture~4 of \citet{meng2026cubefree} states that when $d\mid N$, every projective-$d$-cube-free set $A\subseteq\mathbb Z_N$ satisfies $|A|\le(d-1)N/d$. For each row of Table~\ref{tab:projective-cube-families} and every integer $k\ge1$, let $A_N=\{x\in\mathbb Z_N:x\bmod M\notin F\}$ with $F$ the excluded residues. Reducing a hypothetical cube in $A_N$ modulo $M$ would give a cube that avoids $F$, contradicting the verified obstruction in that smaller group. The corresponding Lean file proves $d\mid N$, that $A_N$ contains no projective $d$-cube, the displayed cardinality, and $d|A_N|>(d-1)N$. Each row therefore refutes Conjecture~4 for infinitely many group orders.

\begin{table}[ht]
\centering
\caption{Three families of counterexamples verified by Lean. Every row
holds for every integer $k\ge1$. The last column compares the density
of the constructed set with the conjectured upper bound.}
\label{tab:projective-cube-families}
\small
\begin{tabular}{@{}rrllll@{}}
\toprule
$d$ & $M$ & Excluded residues $F$ & $N$ & $|A_N|$ & Density \\
\midrule
5 & 16 & $\{0,4,12\}$ & $80k$ & $65k$ & $13/16>4/5$ \\
7 & 24 & $\{0,6,18\}$ & $168k$ & $147k$ & $7/8>6/7$ \\
8 & 36 & $\{0,6,12,30\}$ & $72k$ & $64k$ & $8/9>7/8$ \\
\bottomrule
\end{tabular}
\end{table}

The finite obstruction is checked in four steps. A Python generator enumerates the sets of subset sums reachable while all nonempty sums avoid $F$ modulo $M$; these sets describe the possible states after each new generator is added. Lean checks the generated data: appending any residue must either produce a forbidden nonempty subset sum or lead to a listed state at the next depth, and no state survives to depth $d$. A shared soundness proof turns these checks into a theorem about every length-$d$ sequence. Reduction modulo $M$ then gives cube-freeness for every group order divisible by $M$, counting the repeated residue classes gives the cardinality, and $N=\operatorname{lcm}(M,d)k$ supplies the hypothesis $d\mid N$. The three self-contained files compiled under Lean 4.30.0 and 4.34.0 with Lean's bundled libraries and no Mathlib; the batch checks use kernel evaluation and no axioms, the assembled closure theorem uses \texttt{propext}, and the family theorems use \texttt{propext}, \texttt{Classical.choice}, and \texttt{Quot.sound}. The certificates are in \path{mathematics/theorems/}, with build evidence, generator and template hashes in \path{mathematics/builds/cube_families.json}. A separate model reader examined the argument adversarially, regenerated all three files byte for byte, rejected thirteen deliberately incorrect variants, and reran the $d=5$ and $d=7$ files under both Lean versions; for $d=8$ it inspected the recorded compilation logs and recomputed the mathematics independently.

The $d=5$ set was found in the model runs of 10 September 2026, first with BugBank and then with plain prompting. It is the pullback under reduction modulo 16 of the published layered set $C_5$ of \citet{long2019projective}, so the search rediscovered that construction. We found the $d=7$ and $d=8$ sets by computer search, with AI tools, on 19 September and then formalized them; their construction priority is not established. Human reviews of the dimension-5, 7 and 8 dossiers (R02, R13, R14) are reported in Appendix~\ref{app:recent-human-review}. The same source's auxiliary Conjecture~10 had been refuted publicly on 26 July 2026 by \citet{hypnopump2026conj10} with a different witness and a Lean proof; that witness contains a projective cube, so it concerns Conjecture~10, not Conjecture~4.

The three families raise the question in which dimensions the layered construction itself exceeds the conjectured density. Write $C_d\subseteq\mathbb Z_{2^n}$ for the recursively defined layered construction of \citet{long2019projective}; each layer has an index, and the construction fits in the group when all those indices are at most $n$, which $d\le n$ guarantees. Let $g(d)$ be the fraction of the group outside the construction, so that $g(1)=1$ and $g(P+m-1)=(1+g(m))/(2P)$ for $P=2^\ell$, $\ell\ge1$, and $1\le m\le P$. The Lean proofs establish, for every $d\ge2$ and every $n$ for which the layers fit, that $|C_d|/2^n=1-g(d)$ and that
\[
 \frac{|C_d|}{2^n}>\frac{d-1}{d}
 \quad\Longleftrightarrow\quad
 d\notin\{2^\ell-2,\;2^\ell-1,\;2^\ell:\ell\ge1\};
\]
equality holds exactly when $d$ is a power of two, and at the other exceptional dimensions, one or two below a power of two, the density is strictly below the bound. The theorem \texttt{LongWagner.C\_five\_mod\_sixteen} identifies $C_5$ in $\mathbb Z_{16}$ with the certified quotient excluding $\{0,4,12\}$; at $d=7,8$ the layered densities are $27/32<6/7$ and $7/8$, which is why the other constructions were needed there. Cube-freeness of the general $C_d$ is supplied by \citet[Section~3.1 and Claim~3.1]{long2019projective} and enters the conditional Lean theorems as an explicit premise, while the certificates for $d=5,7,8$ verify cube-freeness directly; a general pullback lemma proves that taking a pullback preserves cube-freeness and density. These results are in \path{mathematics/extensions/layered_boundary/}: \path{LayeredBoundaryAllD.lean} contains \texttt{density\_gt\_iff}, \texttt{density\_eq\_iff}, \texttt{card\_C\_iff\_fits}, \texttt{card\_C\_gt\_iff}, and \texttt{card\_C\_eq\_iff}, and \path{PullbackLemma.lean} the pullback lemmas, under Lean 4.30.0 with Mathlib. The original all-dimensions file has a successful execution receipt; intermediate bounds were then appended, leaving the headline classification unchanged, and a replay of the unedited amended bytes on 21~September 2026 exited with status zero and reproduced the logs byte for byte (\path{LayeredBoundaryAllD_RECEIPT.json}). The scoped supplement L01 reviews this all-dimensional extension with the Long--Wagner cube-freeness premise explicit (Appendix~\ref{app:recent-human-review}).

\paragraph{Supertoken graphs.}
A vertex of the supertoken graph $F_k(K_n)$ is a placement of $k$ indistinguishable tokens among $n$ types, allowing several tokens of the same type \citep[Definition~2.1]{hammacksmith2017reducedpowers}, represented by a list of $n$ natural numbers summing to $k$. Adjacent vertices differ by moving one token to another coordinate. The proof in \path{mathematics/extensions/06_supertoken_metric/SupertokenMetricGeneral.lean} starts from this adjacency and derives the shortest-walk distance as half the $\ell_1$ distance, the formula stated by \citet[Proposition~3.1(i), Eq.~(2)]{baskoro2025supertoken}. A resolving set is a set of landmarks whose distance vectors distinguish all vertices~\citep{harary1976metric,slater1975leaves}. The finite witness supplies six landmarks resolving $F_2(K_8)$; padding each landmark with zero coordinates and adding $2e_i$ for each new coordinate gives at most $n-2$ landmarks for every $n\ge8$. The theorem \texttt{conjecture\_4\_2\_fails\_k2} thus refutes the claimed dimension $n-1$ throughout that range, including the claim that it holds for all sufficiently large $n$ at $k=2$. The result is an upper bound at fixed $k=2$; the exact dimensions, and the behavior as $k$ grows with $n$ fixed, are separate questions, and the proof is independent of the separate expensive lower-bound computation for $(8,2)$. We later extended the family and formalized it with AI tools.

\paragraph{Binomial polynomials.}
Coefficientwise $q$-log-convexity requires every coefficient of $V_{n-1,m}V_{n+1,m}-V_{n,m}^2$ to be nonnegative~\citep{liuwang2007logconvexity}, where $V_{n,m}(q)=\sum_{j=0}^{n}\binom{n+j}{j}^{m}q^j$. The recorded witness has $m=11$ and $n=3$. The coefficient responsible for that failure yields a uniform obstruction: for every $m\ge11$, the coefficient of $q^2$ at $n=3$ is $2\cdot15^m+6^m-2\cdot10^m-16^m<0$, and a kernel-checked calculation at $n=23$ covers the remaining exponents $4\le m\le10$. Together these give \texttt{not\_qlogconvex\_of\_four\_le} in \path{mathematics/extensions/09_sun_qlogconvex/SunQLogConvex.lean}, which refutes the second assertion in the cited arXiv v4 of \citet{sun2025polynomials} for every integer $m\ge4$. The other three assertions, the cases $m=2,3$, and pointwise nonnegativity for real $q>0$ are distinct questions.

\paragraph{Fake-Gaussian positivity.}
For the fixed-$m$ positivity assertion of \citet{gatzweilerkrattenthaler2025positivity}, let
\[
 F(q)=\prod_{i=1}^{21}\left(\frac{1-q^{22+i}}{1-q^i}\right)^{a_i},
 \qquad a_6=a_{21}=1,\quad a_{14}=a_{18}=a_{20}=2,
\]
with all other $a_i=0$. The certificate \path{mathematics/extensions/07_fake_gaussian/FakeGaussian.lean} checks the full numerator--denominator identity and proves that $F$ is the exact integer polynomial represented by its 177 coefficients, with $[q^{43}]F=-1$. The declaration \path{not_conjecture4} refutes the encoded general statement, and a general quotient-uniqueness lemma ensures that every polynomial representative of the quotient has this same negative coefficient. The result concerns Conjecture~4; the symmetric positivity assertion, the source's other conjectures, and eventual positivity as $m$ grows have separate mathematical scope. A model run produced the witness, and we formalized it with AI tools; Appendix~\ref{app:recent-human-review} reports the human review of the dossier (R07) and the scoped supplement G07.

\paragraph{$q$-rious unimodality.}
For the factorial-ratio unimodality assertion of \citet{warnaarzudilin2025qrious}, the recorded tuples are $a=(2,3,5,12)$ and $b=(1,1,1,4,6,9)$. Writing $[j]_q=1+q+\cdots+q^{j-1}$, their quotient satisfies $D(a,b;q)=[10]_q[11]_q(1-q^2+q^4)$, and the full coefficient sequence of $(1+q)D$ is not unimodal. The certificate \path{mathematics/extensions/04_qrious_unimodality/QriousUnimodality.lean} proves the quotient identity and its uniqueness and derives the general negations \path{not_conjecture5} and \path{not_conjecture5Sym}. The Landau hypothesis~\citep[Eq.~(1)]{warnaarzudilin2025qrious} is certified for every general floor profile, including every rational input, and for every structure satisfying the listed floor properties; the written argument identifies the real numbers with such a structure, while the retained Lean proof covers the abstract structure and rational inputs. The refutation concerns unimodality and leaves the separate positivity conjecture as a distinct question. The witness came from a model run; the human review of the dossier (R04) is reported in Appendix~\ref{app:recent-human-review}.

\paragraph{Permutation-grid rooks.}
For the permutation $w=(2,7,4,8,1,6,3,5)$, the Lean kernel computes 155 complete non-attacking rook placements on its permutation crossword grid; since $155\equiv3\pmod4$, this count violates the proposed condition. Using general definitions of permutation grids, words, and complete placements, the theorem \texttt{not\_conjecture310} in \path{mathematics/supporting/sources/RookSpectrum.lean} refutes the necessity direction of Conjecture~3.10 of \citet{lewiswon2026permutationgrids}. Two size ranges matter here: the board has size eight, within the board sizes the source enumerated, whereas the count 155 exceeds the $r\le108$ count range summarized immediately before the conjecture, and the restatement given to the model omitted that source sentence. The certificate establishes this witness and this direction of the conjecture; the converse direction, a corrected spectrum of possible counts, and the smallest possible witness are further questions. The human review of the dossier (R05) is reported in Appendix~\ref{app:recent-human-review}.

\paragraph{Toda numerators.}
A coefficient sequence is unimodal if it never decreases before a peak and never increases after it. The recorded Toda witness of 10 September has numerator $1+q^2$: its coefficients are not unimodal when consecutive powers of $q$ are counted, but become unimodal when the polynomial is written in $q^2$, and under that substitution it is the source's type-$A_2$ example. This suggested a stronger target, a failure that persists when the polynomial is written in any fixed integer power of $q$, a change called common-exponent regrading. While formalizing the result on 19 September with AI tools, we found a type-$B_2$ witness at $\alpha=2\alpha_1+2\alpha_2$ with numerator $1+q+3q^2+2q^3+5q^4+2q^5+3q^6+q^7+q^8$. Its nonzero coefficient of $q$ forces the regrading step to be one, and $3>2<5$ gives an interior dip. The theorem \texttt{Exact.not\_conjecture73\_exact} in \path{mathematics/supporting/sources/TodaNumerator_Exact.lean} formalizes the refutation for general root-system data using the exact defining recurrence after denominators are cleared, checking full polynomial identities. The connection from that recurrence to the rational functions in Definition~1.1 of \citet{labelle2026toda} uses algebra in $\mathbb Q(q)$ and the source's uniqueness result in the written argument. The original model witness and the later, stronger witness retain separate discovery records; the human review of the dossier (T11) is reported in Appendix~\ref{app:recent-human-review}.

\paragraph{Checking definitions beyond the witness.}
A witness can satisfy two definitions that disagree elsewhere. For the $C_2^4$ short-zero-sum record, replacing the last invariant factor by the first in the definition of the exponent, reversing divisibility, or using modulus $2$ throughout still lets the refutation compile, because all of these agree at $[2,2,2,2]$. A companion file gives a separately written, Lean-equivalent formulation and checks true instances at $C_2\oplus C_4$, which the altered definitions fail. These mutation probes~\citep{jia2011mutation} test the definitions beyond the refuting instance; correspondence to the source prose is reviewed separately. The base and companion proofs, with their successful compilation logs, are listed in \path{mathematics/supporting/index.json} as \texttt{ShortZeroSums\_Conjecture.lean} and \texttt{ShortZeroSums\_Conjecture\_Sanity.lean}; the individual mutant logs were not retained.

\paragraph{What was contributed at each stage.}
The research collection contains eleven model-produced refutation cases and one finding already visible within its source. Subsequent checking identified known constructions, including the cube and Schur examples, an index-ratio conjecture that its author had revised before our run, and a semigroup contradiction present in the source's own table. We keep these distinctions so that a reader can see both the mathematical conclusion and what was contributed at each stage: identifying a witness, checking it against the source, extending its range, or formalizing the result. The cases are reported as source-linked mathematical criticism and formal completion, with discovery attribution stated individually.
